% soundness-completeness.tex

%%%%%%%%%%%%%%%%%%%%
\begin{frame}{}
  \begin{center}
    \purple{\Large 就是现在, 将语法与语义统一起来!!!}

		\vspace{-0.30cm}
		\[
			\blue{\Sigma \models \alpha}
		\]
		\vspace{-0.50cm}
    \fig{width = 0.50\textwidth}{figs/syntax-semantics}
		\[
			\brown{\Sigma \vdash \alpha}
		\]

    \pause
		\vspace{-1em}
		\[
			\text{如果}\; \brown{\Sigma \vdash \alpha},
			\;\text{则}\; \blue{\Sigma \models \alpha}
			\pause\qquad\qquad
			\text{如果}\; \blue{\Sigma \models \alpha},
			\;\text{则}\; \brown{\Sigma \vdash \alpha}
		\]
  \end{center}
\end{frame}
%%%%%%%%%%%%%%%%%%%%

%%%%%%%%%%%%%%%%%%%%
\begin{frame}{}
  \begin{theorem}[命题逻辑的可靠性 (Soundness)]
    如果 $\Sigma \vdash \alpha$, 则 $\Sigma \models \alpha$。
  \end{theorem}

  \vspace{0.30cm}
  \fig{width = 0.50\textwidth}{figs/syntax-semantics}
  \vspace{0.30cm}

  \begin{theorem}[命题逻辑的完备性 (Completeness)]
    如果 $\Sigma \models \alpha$, 则 $\Sigma \vdash \alpha$。
  \end{theorem}
\end{frame}
%%%%%%%%%%%%%%%%%%%%